\documentclass[11pt]{article}
\usepackage[margin=1in]{geometry}
\usepackage{hyperref}
\usepackage{enumitem}
\title{Phase 2 Technical Report: Task-Specific Meta-Path Compilation}
\author{RFM v2 Phase 2}
\date{\today}

\begin{document}
\maketitle

\section{Objective}
Phase 2 consumes Phase 1 artifacts and a task definition, then compiles a task-specific policy: selected meta-paths, attribute priorities, aggregation hints, and operational budgets.

\section{Inputs and Outputs}
\textbf{Inputs}: dataset name, task name, Phase 1 artifact directory, LLM model.

\textbf{Outputs} under \texttt{<output-dir>/<dataset>/<task>/}:
\begin{itemize}[leftmargin=1.5em]
  \item \texttt{task\_spec.json}
  \item \texttt{final\_validation\_report.json}
  \item \texttt{solver\_rounds.jsonl}
  \item \texttt{critic\_rounds.jsonl}
\end{itemize}

\section{Refactored Module Layout}
\begin{itemize}[leftmargin=1.5em]
  \item \texttt{phase2\_pipeline.py}: orchestration, context assembly, solver-critic loop.
  \item \texttt{phase2\_policy.py}: selection sanitization, finalization, aggregation-plan build, policy merge.
  \item \texttt{phase2\_rescoring.py}: LLM-based candidate path rescoring and fallback parsing.
  \item \texttt{phase2\_prompts.py}: solver and critic prompt templates.
  \item \texttt{phase2\_llm.py}: local JSON-oriented chat generation wrapper.
  \item \texttt{phase2\_io.py}, \texttt{phase2\_validation.py}: artifact IO and schema/consistency checks.
\end{itemize}

\section{Compilation Procedure}
\begin{enumerate}[leftmargin=1.8em]
  \item Load Phase 1 artifacts including \texttt{semantic\_context\_graph.json}.
  \item Build candidate meta-paths from anchor table over semantic graph edges.
  \item Score candidate paths with an LLM-based relevance pass.
  \item Run solver-critic iterations to refine selected paths and policy fields.
  \item Apply deterministic finalization (anchor/depth/diversity constraints).
  \item Emit final \texttt{task\_spec.json} and validation report.
\end{enumerate}

\section{Rescoring Design}
The rescoring module supports robust JSON variability:
\begin{itemize}[leftmargin=1.5em]
  \item accepts multiple output shapes (list/map/alternate keys);
  \item path-id normalization (e.g., \texttt{p1}, \texttt{P0001}, embedded IDs);
  \item order-based fallback when IDs are malformed or omitted;
  \item compact second-pass fallback returning a numeric relevance list.
\end{itemize}

\section{Selection Policy}
Current selection behavior combines:
\begin{itemize}[leftmargin=1.5em]
  \item LLM task relevance,
  \item structural priors (depth/cardinality),
  \item directness and dominated-path suppression,
  \item safety defaults (exclude low-confidence links by default).
\end{itemize}

\section{Diagnostics}
Console diagnostics expose:
\begin{itemize}[leftmargin=1.5em]
  \item semantic graph usage summary,
  \item SQL hint coverage,
  \item LLM rescoring stats (scored/matched/returned/batch failures),
  \item solver-critic per-round status and issue summaries,
  \item final top selected paths with relevance scores.
\end{itemize}

\section{Extensibility}
\begin{itemize}[leftmargin=1.5em]
  \item Replace/extend path rescoring in \texttt{phase2\_rescoring.py}.
  \item Swap selection/finalization rules in \texttt{phase2\_policy.py}.
  \item Add prompt constraints in \texttt{phase2\_prompts.py} without changing orchestration.
  \item Add validation invariants in \texttt{phase2\_validation.py}.
\end{itemize}

\section{Reproducibility}
\begin{verbatim}
python phase2/phase2_pipeline.py \
  --dataset rel-f1 \
  --task driver-position \
  --phase1-artifacts-dir phase1/artifacts \
  --output-dir phase2/artifacts \
  --model-size 8b \
  --max-rounds 5
\end{verbatim}

\end{document}
